\section{The cubic branch filling and the normalized cubic identity}
\label{sec:normalized-cubic-identity}

We use the standard full genus-zero GT action and the exact
nonabelian $(K_3)$ identity. The rigidity input is
\begin{equation}\label{eq:cubic-normalizer-1}
\operatorname{Out}(\widehat\Gamma_{0,n})
 =\widehat{\mathrm{GT}}\times\mathfrak S_n\qquad(n\ge5),
\end{equation}
as proved by Hoshi--Minamide--Mochizuki
\cite[Definition~2.7 and Corollary~2.8(i), author pp.~31--32]{CubicHMM}.
Here the set of primes is the full set; their Definition~2.7
identifies the resulting GT group with the classical full
profinite GT group. The other topological inputs and their
precise scope are stated where they are used.

Schneps described the cyclic genus-zero quotient construction and its
discrete normalizers in \cite[Section~4.3,
Theorems~4.3.1--4.3.2, author pp.~54--58]{SchnepsSpecialLoci}.
The mixed version and its full profinite topology are checked below.
She introduced special homomorphisms from these cyclic automorphisms
of punctured spheres. Her subgroup $\widehat{\mathrm{GS}}^1_{0,0}$
is defined by two such compatibilities, and its elements satisfy
the normalized cubic relation $(\mathrm{NT}_2)$
\cite[Definition~1.2, p.~658; Theorems~3.1 and~3.4,
pp.~665 and~667]{Schneps2006Special}.
The argument here proves the cubic comparison on full GT with
multiplier one and vanishing $\rho_3^I$.


Throughout, $\Gamma_{0,n}$ denotes the pure mapping class
group of an oriented sphere with $n$ marked points. Brackets
around a set of marked points mean that its permutations are
allowed. A circumflex denotes full profinite completion.
Let $g=(1,f_g)$ belong to full profinite GT. Its standard genus-zero
action \cite[equation~(1.3), p.~657]{Schneps2006Special}
is represented by an automorphism $\alpha_g$
of each completed mixed mapping class group. It induces the
identity on the finite permutation quotient and the identity
on the abelianization of each pure subgroup. We distinguish
this action from its outer class on a pure subgroup.
The labelled forgetting maps are compatible with these actions
up to inner automorphisms by pure elements. One may choose
compatible based representatives for a single forgetting map
and obtain equality; changing such a choice adds a pure inner
automorphism. This is the usual labelled genus-zero action.
Below we keep these
pure inner corrections explicitly. On the four-point group it
has the usual outer class of
\begin{equation}\label{eq:cubic-normalizer-2}
\phi_g(x)=x,\qquad
 \phi_g(y)=f_g(x,y)^{-1}y f_g(x,y).
\end{equation}
The multiplier-one hypothesis is used here.

\subsection{The centralizer and its full completion}

Fix a primitive cube root $\omega$. On $\mathbb P^1$ mark
$0,\infty$ and three free orbits
\[
 D_{i,j}=\omega^j t_i,
 \qquad i\in\{1,2,3\},\quad j\in\mathbb Z/3,
 \qquad t_i\ne0,\quad t_i^3\ne t_k^3\ (i\ne k).
\]
Thus there are eleven marked points. Let $h(z)=\omega z$,
$\mathcal H=\langle h\rangle$, and let $\mathcal G$ be the subgroup of the
mixed mapping class group whose permutations fix $0,\infty$
individually and commute with the displayed permutation of
$h$. Put $\mathcal N=C_{\mathcal G}(\mathcal H)$. The cyclic quotient has coordinate
$w=z^3$, with its two branch points and three further marked
points. We first prove the discrete centralizer sequence
\begin{equation}\label{eq:cubic-normalizer-3}
1\longrightarrow \mathcal H\longrightarrow \mathcal N\longrightarrow
 Q_3\longrightarrow1,
 \qquad Q_3=\Gamma_{0,\{0\},\{\infty\},[3]}.
\end{equation}
The argument works with two or three free orbits. Delete all
marked points from the covering sphere and its quotient.
The resulting cover $S\to S/\mathcal H$ is unbranched, regular and
three-sheeted. Its monodromy character takes the values $1$
and $-1$ on positive loops around the two distinguished
branch labels and zero on the other punctures.

Every mapping class centralizing $h$ has a
representative commuting with $h$. Give $S$ its displayed
complex structure, in which $h$ is conformal, and choose the
extremal quasiconformal representative $n$ of that class.
Existence and uniqueness for a finitely punctured compact
surface follow from Gardiner--Hu, Section~5, Theorem~1,
and Section~8, Theorem~6
\cite[author pp.~14--15 and~22]{CubicGH}; the punctured
hypothesis is specified immediately before their Theorem~1.
Both $nh$ and $hn$ lie in the same homotopy class, and both
are extremal: any improvement, composed with $h^{-1}$ on
the appropriate side, would improve $n$. Uniqueness implies
that they differ by a conformal map homotopic to the identity.
Such a map fixes all punctures, hence at least three points
of the compactified sphere, and is the identity. Thus $nh=hn$.
This spells out in the punctured setting the extremal-map
argument of Birman--Hilden \cite[p.~433]{CubicBH}.

The representative $n$ descends to the punctured quotient.
If two such representatives are isotopic upstairs, their
quotient is fiber-preserving and isotopic to the identity.
Birman--Hilden Theorem~1
\cite[p.~424]{CubicBH}, applied to this unbranched cover,
makes it fiber-isotopic to the identity. Their descended
mapping classes consequently agree. Conversely every element
of $Q_3$ preserves the monodromy character and therefore
lifts; every lift commutes with $h$ and has exactly the
required label permutation. The lifts differ by $\mathcal H$.
Lifting an isotopy downstairs shows that the kernel is $\mathcal H$,
which proves \eqref{eq:cubic-normalizer-3}. This argument proves the required
liftability directly. It does not use Birman--Hilden's later
general lifting assertion in Lemma~5.1 or Theorem~5;
their erratum leaves Theorems~1--4 unaffected
\cite{CubicBHErr}.

We prove that $\mathcal N$ has its full profinite topology in the
ambient mapping class group. Let
$\Phi=\pi_1(S/\mathcal H)$ and $\Pi=\pi_1(S)$, with compatible
basepoints. Thus $\Pi\lhd\Phi$ is the kernel of the above
character $\chi:\Phi\to C_3$, and
$\widehat\Pi$ is open in the free profinite group
$\widehat\Phi$. Put
\[
 \mathcal L=
 \operatorname{Aut}_{\chi}(\widehat\Phi)
                 /\operatorname{Inn}(\widehat\Pi),
\]
where the subscript means preservation of $\chi$ pointwise.
Since $\widehat\Phi$ is finitely generated and
$\operatorname{Inn}(\widehat\Pi)$ is compact, these are
profinite groups. There is an exact sequence
\begin{equation}\label{eq:cubic-normalizer-4a}
1\longrightarrow C_3\longrightarrow\mathcal L
 \longrightarrow\operatorname{Out}_{\chi}(\widehat\Phi)
 \longrightarrow1.
\end{equation}
Indeed, its kernel is
$\operatorname{Inn}(\widehat\Phi)/
\operatorname{Inn}(\widehat\Pi)=\widehat\Phi/\widehat\Pi$;
the centre of a nonabelian free profinite group is trivial.

The lifted mapping classes give a homomorphism
$\mathcal N\to\mathcal L$. Choose a path upstairs from the basepoint
to its image under a commuting representative and project
that path downstairs. The resulting based automorphism of
$\Phi$ is well-defined modulo $\operatorname{Inn}(\Pi)$:
changing the path gives just such an inner automorphism, and
changing the representative is harmless by the preceding
fiber-isotopy argument. This construction is compatible with
composition. Its profinite extension fits with \eqref{eq:cubic-normalizer-3}, \eqref{eq:cubic-normalizer-4a}
and the ordinary outer action of $\widehat Q_3$ on
$\widehat\Phi$.

The completion of \eqref{eq:cubic-normalizer-3} has kernel exactly $\mathcal H$. This finite
subgroup is separated by its permutation of the upstairs
punctures; the universal property of completion identifies
$\widehat{\mathcal N}/\mathcal H$ with $\widehat Q_3$. Asada's faithfulness
theorem makes
$\widehat Q_3\to\operatorname{Out}(\widehat\Phi)$
injective in the pure case
\cite[introduction Theorem, p.~3, parts (i)--(ii)]{CubicAsada}.
The mixed case follows as well: a nontrivial puncture
permutation is visible on peripheral conjugacy classes,
and its permutation kernel is the pure subgroup. The map
$\mathcal H\to C_3$ in \eqref{eq:cubic-normalizer-4a} is an isomorphism, since a deck
transformation is represented by a projected path of the
corresponding monodromy. The two exact sequences therefore
make $\widehat{\mathcal N}\to\mathcal L$ injective.

Restriction $\mathcal L\to
\operatorname{Out}(\widehat\Pi)$ is injective. If a
representative restricts to an inner automorphism, adjust by
that inner automorphism so that it fixes $\widehat\Pi$
pointwise. For any $x\in\widehat\Phi$, comparison of its
conjugation on $\widehat\Pi$ then shows that
$x^{-1}\theta(x)$ centralizes $\widehat\Pi$.
This centralizer is trivial: it intersects $\widehat\Pi$
in the latter's trivial centre, so injects into the finite
group $\widehat\Phi/\widehat\Pi$; a free profinite group
has no nontrivial finite subgroup. Thus $\theta$ is the
identity. We have proved faithfulness of the full completed
$\mathcal N$-action on $\widehat\Pi$. This action factors through
the completion of the ambient mapping class group, so
$\widehat{\mathcal N}\to\widehat{\mathcal G}$ is injective. In particular
\begin{equation}\label{eq:cubic-normalizer-4}
\overline{\mathcal N}=\widehat{\mathcal N},\qquad
 \overline{\mathcal N}/\mathcal H=\widehat Q_3.
\end{equation}
This establishes the full profinite topology on $\mathcal N$.

Finally the Boggi--Zalesskii centralizer theorem gives
$C_{\widehat{\mathcal G}}(\mathcal H)=\overline{\mathcal N}$
\cite[Theorem~A(iii), p.~3; Section~6.1 and
Theorem~6.2(iii), pp.~16--17]{CubicBZ}.
The theorem applies to the present finite-index subgroup:
genus-zero mapping class groups are good and have finite
virtual cohomological type by the cited Section~6.1;
goodness and that finiteness property pass to finite-index
subgroups (by the finite-index cohomology comparison and
Shapiro's lemma). Residual finiteness is inherited as well.
Only finite $3$-subgroups are used.

\subsection{Correction and identification of the parameter}

Let $P=\Gamma_{0,11}$ and let
$\mathcal G_{\mathrm{cyc}}=P \mathcal H$ be the inverse image of the cyclic permutation
subgroup generated by $h$. Since $\alpha_g$ is the identity
on label permutations, $\alpha_g(\mathcal H)\subset\widehat{\mathcal G}_{\mathrm{cyc}}$.
The Boggi--Zalesskii finite-$3$-subgroup theorem
\cite[Theorem~A(i), p.~3; Theorem~6.2(i), p.~17]{CubicBZ}
applied to $\mathcal G_{\mathrm{cyc}}$ conjugates $\alpha_g(\mathcal H)$ into $\mathcal G_{\mathrm{cyc}}$.
The conjugator can be taken pure: write an initial conjugator
as $h^j p$ with $p\in\widehat P$, and absorb $h^j$ by
conjugating the resulting discrete subgroup. The resulting
generator has the same permutation of all eleven labels as
$h$, since the permutation quotient of $\mathcal G_{\mathrm{cyc}}$ is abelian.

Nielsen realization realizes this finite cyclic group by
conformal transformations of a punctured sphere. The required
finite-type, punctured scope is included in Kerckhoff's
extension to finitely generated Fuchsian groups with cusps
\cite[Section~4, pp.~255--256]{CubicKerckhoff}.
Putting its fixed points at $0,\infty$, the generator is
$z\mapsto\omega z$ or $z\mapsto\omega^{-1}z$.
In the first case its labelled configurations form the
connected space
\[
 \{(t_1,t_2,t_3)\in(\mathbb C^\times)^3:
                         t_i^3\ne t_k^3\text{ for }i\ne k\}.
\]
In the second case exchange the two fixed-point labels first.
It follows that there is an $e$ which is pure, or pure
followed by a transposition of $0,\infty$, and which fixes
every free-orbit label, such that
\begin{equation}\label{eq:cubic-normalizer-5}
\alpha'=\operatorname{Ad}(e^{-1})\alpha_g
       \quad\text{fixes }\mathcal H\text{ pointwise}.
\end{equation}
Here the assertion about the labels concerns the permutation
of $e$; its pure factor need not fix points pointwise as a
homeomorphism.

We verify invariance of the mixed groups before taking
outer classes. The permutation of $e$ is either the identity
or the single transposition $(0\ \infty)$; it fixes all nine
free labels. Either permutation centralizes the permutation
subgroup defining $\mathcal G$. Since $\alpha_g$ induces the identity
on permutations, $\alpha'$ preserves $\widehat{\mathcal G}$ and
induces the identity on its permutation quotient. Since it
fixes $\mathcal H$ pointwise, it preserves
$C_{\widehat{\mathcal G}}(\mathcal H)=\widehat{\mathcal N}$. Equation \eqref{eq:cubic-normalizer-4} therefore
gives an automorphism $\delta_3$ of $\widehat Q_3$.
It preserves its pure subgroup
$P_5=\widehat\Gamma_{0,5}$ and induces the identity on
$\widehat Q_3/P_5=\mathfrak S_3$.

Write $\vartheta=[\delta_3|_{P_5}]=(g',\sigma)$ by \eqref{eq:cubic-normalizer-1}.
For each $\pi\in\mathfrak S_3$, choose $z_\pi\in
\widehat Q_3$ over $\pi$. Conjugation by $z_\pi$ on
$P_5$ has outer class $(1,\pi)$ in \eqref{eq:cubic-normalizer-1}. As
$\delta_3(z_\pi)$ has the same permutation, the identity
\[
 \delta_3\operatorname{Ad}(z_\pi)\delta_3^{-1}
       =\operatorname{Ad}(\delta_3(z_\pi))
\]
restricted to $P_5$ gives
$\vartheta(1,\pi)\vartheta^{-1}=(1,\pi)$.
Hence $\sigma\pi\sigma^{-1}=\pi$ for every $\pi$.
The centralizer in $\mathfrak S_5$ of the full symmetric
group on the three orbit labels consists exactly of the
identity and the transposition of the other two labels.
Thus
\begin{equation}\label{eq:cubic-normalizer-6}
\sigma=1\quad\text{or}\quad\sigma=(0\ \infty).
\end{equation}
To forget the third orbit, first restrict $\widehat{\mathcal N}$
to the subgroup $\widehat{\mathcal N}^{\rm orb}$ fixing each free
orbit as a set; its quotient by $\mathcal H$ is $P_5$. This subgroup
is preserved by $\alpha'$, as just proved. Ambient forgetting
of all three labels in the third orbit maps it onto the
two-orbit centralizer $\widehat{\mathcal N}_2^{\rm orb}$.
On quotients by $\mathcal H$ this is the usual surjection
$P_5\to P_4=\widehat\Gamma_{0,4}$, and it is the
identity on $\mathcal H$. These facts prove surjectivity upstairs too.

Here are the inner corrections in this step. Choose standard
GT and permutation representatives on $P_5$. By definition
of an outer class, $\delta_3|_{P_5}$ is their product
followed by $\operatorname{Ad}(a_3)$ with $a_3\in P_5$.
Both possible $\sigma$ fix the deleted orbit label.
The standard representatives therefore preserve the
forgetting kernel, and so does $\delta_3$; the correction
becomes the image of $a_3$ in $P_4$. Thus the
induced automorphism $\delta_2$ has outer class
$(g',\sigma)$ on $P_4$.
On the ambient mixed group the same conclusion follows from
labelled forgetting compatibility. If a choice of based
representatives introduces a pure correction $c$, the
forgotten action is
$\operatorname{Ad}(\bar e^{-1}c)\alpha_g$.
Its effective conjugator $c^{-1}\bar e$ still fixes every
remaining free label and has exactly the same possible
branch transposition as $e$. Thus the label and homology
arguments below are unchanged. This notation also covers
all later forgetting maps: every extra based correction is
pure in its own target.

We make the resulting two-orbit calculation explicit. Scale
the first orbit to $\{1,\omega,\omega^2\}$, write the
second as $\{t,\omega t,\omega^2t\}$, and put $w=t^3$.
Then
\[
 \mathcal N_2/\mathcal H\cong\pi_1(\mathbb P^1_w\setminus\{0,1,\infty\})
       =F(x,y).
\]
The subgroup $\mathcal N_2^0$ fixing the chosen representatives $1,t$
fixes all six free-orbit labels and intersects $\mathcal H$ trivially.
Its image is
\begin{equation}\label{eq:cubic-normalizer-7}
H_3=\ker(F(x,y)\longrightarrow C_3),\qquad x\mapsto1,
 \quad y\mapsto0.
\end{equation}
It has basis $X=x^3$, $y_j=x^jyx^{-j}$, $j=0,1,2$.
The element $y_j$ is the positive simultaneous collision
multitwist pairing the three points of one free orbit with
the three points of the other, with relative offset $j$.
It is pure. Its abelian class $v_j$ in the ambient pure
eight-point group is fixed by $\alpha'$: $\alpha_g$ acts
as the identity on this abelianization and the only possible
permutation of $e$ exchanges $0,\infty$, fixing all the
paired labels.

The three $v_j$ are distinguished integrally by forgetting
maps. Retain $0,\infty$ and just one pair with offset $j$.
The class $v_j$ maps to the positive generator around their
collision in the four-point group, whereas the other two
classes map to zero. In particular it is nonzero and
primitive in this test quotient.

Put $\lambda'=\lambda(g')$. In \eqref{eq:cubic-normalizer-7}, an automorphism of
outer class $(g',\sigma)$ takes the conjugacy class of $y_j$
to that of
\begin{equation}\label{eq:cubic-normalizer-8}
y_{\varepsilon\lambda'j+c}^{\lambda'},
 \qquad \varepsilon=1\text{ if }\sigma=1,
 \quad\varepsilon=-1\text{ otherwise},
\end{equation}
with subscripts modulo three and a common offset $c$.
For the GT factor this follows directly from
$x\mapsto x^{\lambda'}$, $y\mapsto f_{g'}^{-1}y^{\lambda'}f_{g'}$;
the interchange of $0,\infty$ lifts to $t\mapsto1/t$ and
reflects the three collision points, and an inner correction
adds a common offset. Once the subscript in \eqref{eq:cubic-normalizer-8} has been
chosen, its conjugating element is in $\widehat H_3$, hence
is pure upstairs.
The preceding three abelian tests therefore give
\begin{equation}\label{eq:cubic-normalizer-9}
v_j=\lambda'v_{\varepsilon\lambda'j+c}
 \quad(j=0,1,2),\qquad
 \lambda'=1,\quad c=0,\quad\varepsilon=1.
\end{equation}
Thus $\sigma=1$. Moreover the two-orbit quotient
automorphism has the form
\begin{equation}\label{eq:cubic-normalizer-10}
\delta_2=\operatorname{Ad}(a)\phi_{g'},
                  \qquad a\in\widehat H_3,
\end{equation}
because its common offset is now the image of $a$ in $C_3$.

Forget all but $0,\infty,1,t$ upstairs. The resulting map
on \eqref{eq:cubic-normalizer-7} is exactly the standard cyclic-cover filling
\begin{equation}\label{eq:cubic-normalizer-11}
d_3:\widehat H_3\longrightarrow\widehat F(x,y),
 \qquad X\mapsto x,\quad y_0\mapsto y,
 \quad y_1,y_2\mapsto1.
\end{equation}
The exact $(K_3)$ identity says
$d_3(f_{g'})=y^{c_3(g')}f_{g'}$. In the automorphism formula
this extra power commutes with $y$, so
\begin{equation}\label{eq:cubic-normalizer-12}
d_3\phi_{g'}|_{\widehat H_3}=\phi_{g'}d_3.
\end{equation}
For clarity, \eqref{eq:cubic-normalizer-12} is checked on the four free generators:
on $X$ both sides give $x$; on $y_0$ both give
$f_{g'}^{-1}yf_{g'}$; and on $y_j$ for $j=1,2$ use
\[
 \phi_{g'}(y_j)=(x^j f_{g'}^{-1}x^{-j})\,y_j\,
                         (x^j f_{g'}x^{-j}),
\]
whose image under $d_3$ is one. The conjugating factors
belong to $H_3$ because $f_{g'}$ is a commutator.
Equation \eqref{eq:cubic-normalizer-10} now contributes the inner automorphism
$\operatorname{Ad}(d_3(a))$ in the target. Thus
\eqref{eq:cubic-normalizer-10}--\eqref{eq:cubic-normalizer-12} make the induced outer action downstairs
equal to that of $g'$. On the other hand, the ambient
forgetting compatibility and \eqref{eq:cubic-normalizer-5} make it the action of $g$,
possibly composed with the transposition of the chosen
labels $0,\infty$. That permutation is a single transposition on the
four chosen labels: the two free labels are individually
fixed. It is not a double transposition in the Klein four
kernel of the four-point permutation action. Equivalently,
in the coordinate $t$ it is $t\mapsto1/t$, which exchanges
the peripheral classes at zero and infinity. Both GT
factors and all pure inner corrections preserve the three
peripheral classes individually. The branch transposition
is therefore absent. Thus $e$ is pure and $g'=g$.

For completeness, the usual GT embedding into the outer
automorphism group of $\widehat F(x,y)$ is injective in this
comparison. An inner conjugation equating two standard
representatives with the same multiplier must centralize
$x$, so is a power of $x$. Comparing the conjugators of
$y$ and then taking their exponent sums makes that power
and the possible power of $y$ zero; the two commutator
parameters are equal.

\begin{lemma}[The other cubic filling]\label{lem:cubic-branch-filling}
Let $g=(1,f_g)\in\widehat{\mathrm{GT}}$.
Fill the punctures $0,\infty$ in the threefold cyclic cover
$\mathbb P^1_t\setminus\{0,\infty,1,\omega,\omega^2\}$.
The restriction of the standard action $\phi_g$ on the
base group, followed by this filling, induces on
$\pi_1(\mathbb P^1_t\setminus\{1,\omega,\omega^2\})$
the outer class of the same $g$.
\end{lemma}
\begin{proof}
Retain the three points $1,\omega,\omega^2$ of the first
orbit and the chosen point $t$ of the second, forgetting
$0,\infty$ and the other two points of the second orbit.
On $\mathcal N_2^0=H_3$ this is precisely the displayed branch
filling. The parameter comparison above gives
$\delta_2=\operatorname{Ad}(a)\phi_g$ with
$a\in\widehat H_3$, and the correction $e$ is pure.
Both corrections therefore become inner automorphisms in
the target four-point group. Ambient forgetting supplies
the standard outer action of $g$ there, proving the claim.
\end{proof}

\subsection{The exact cubic conjugator}

Let $\Gamma=\widehat{C_2*C_3}$ be the modular quotient of
the signed braid group on $u,v$, and put $r=uv$,
$a=u^2$, $b=v^2$. The subgroup $\mathcal E=\langle r,a\rangle$
has presentation
\[
 \mathcal E=\langle r,a\mid r^3=(ra)^3=1\rangle
       \cong C_3\amalg C_3.
\]
Its index-three subgroup $F=\langle a,b\rangle$ is free;
conjugation by $r$ cyclically permutes $a,b,(ab)^{-1}$.
It is the filled cover in Lemma~\ref{lem:cubic-branch-filling},
with $x\mapsto r$, $y\mapsto a$. These choices fix the
orientation and the order of the three peripheral loops.
The signs can also be seen directly. Send the distinguished
point $t$ of $\{1,\omega,\omega^2,t\}$ to infinity and make an
affine change of coordinates. The three moving points become
\[
 Z_i(t)=\omega^{2i}+t\omega^i\quad(0\leq i<3),
 \qquad Z^3-3tZ-(1+t^3)=0.
\]
A positive loop of $w=t^3$ around zero ends at $\omega t$;
the endpoint configuration is $\omega^2$ times the initial
one. Closing by a positive one-third rotation gives $uv=r$.
At $w=1$, the difference
$Z_1(t)-Z_2(t)=(\omega^2-\omega)(1-t)$ makes a positive
full turn, giving $u^2=a$. A full-rotation ambiguity is central
and disappears in the modular quotient.

\begin{proposition}\label{prop:normalized-HF1-from-normalizer}
If $g=(1,f_g)\in\widehat{\mathrm{GT}}$ and
$\rho_3^I(g)=0$, then its pure cubic conjugator satisfies
\begin{equation}\label{eq:cubic-normalizer-13}
\boxed{\quad A_g(u^2,v^2)^{-1}=f_g(u^2,uv)\quad}
\end{equation}
in the full signed profinite braid group.
\end{proposition}
\begin{proof}
The standard action $\phi_g$ descends to $\mathcal E$: it preserves
the two finite peripheral relations because it sends each
peripheral generator to a conjugate of itself. Denote this
action by $\psi$. Put $H_g=f_g(a,r)$. The two-cycle relation
gives
\begin{equation}\label{eq:cubic-normalizer-14}
\psi(r)=r,\qquad\psi(a)=H_g aH_g^{-1}.
\end{equation}
The braid action $\Phi_g$ preserves $\mathcal E$ and has
$\Phi_g(a)=a$, $\Phi_g(r)=A_g rA_g^{-1}$.
Lemma~\ref{lem:cubic-branch-filling} says that on the open
free subgroup $F$ one has
\begin{equation}\label{eq:cubic-normalizer-15}
\psi|_F=\operatorname{Ad}(k)\Phi_g|_F
                  \quad\text{for some }k\in F.
\end{equation}
The inner correction is in $F$, since the target four-point
group is pure and all corrections in the preceding lemma
were pure. Two automorphisms of $\mathcal E$ agreeing on $F$ agree
everywhere: their quotient fixes $F$ pointwise, and comparison
of conjugation on $F$ shows that its discrepancy at every
element centralizes $F$, whose centralizer is trivial.
Thus \eqref{eq:cubic-normalizer-15} holds on $\mathcal E$.

Applying it to $r$ gives $kA_g\in C_\mathcal E(r)=\langle r\rangle$.
Since $k,A_g\in F$ and $F\cap\langle r\rangle=1$, this
forces $k=A_g^{-1}$. Applying \eqref{eq:cubic-normalizer-15} to $a$ now gives
$A_gH_g\in C_F(a)=\overline{\langle a\rangle}$.
Write $A_gH_g=a^t$ with $t\in\widehat{\mathbb Z}$.

The factor $A_g$ has abelianization zero. This follows from
$f_g=\operatorname{Ad}(r)(A_g)A_g^{-1}$: on $F^{\rm ab}$
the order-three rotation minus the identity has determinant
three, and multiplication by three on $\widehat{\mathbb Z}$
is injective, so that rotation has no fixed vector.
Write the class of $f_g(x,y)$ in the full Schreier
abelianization of $H_3$ as
\[
 (0,\psi_0,\psi_1,\psi_2),\qquad
 \psi_0+\psi_1+\psi_2=0.
\]
The first coordinate and the displayed sum vanish because
$f_g$ is a commutator. Ihara's full positive-root normalization
is $\rho_3^I=-\psi_0$
\cite[Section~5, p.~205]{NT2002}.
Branch filling sends $y_0,y_1,y_2$ to $a,b,(ab)^{-1}$.
The two-cycle relation therefore gives
\[
 [H_g]=(-\psi_0+\psi_2,-\psi_1+\psi_2)=(t,0).
\]
The second coordinate gives $\psi_1=\psi_2$, while
$\rho_3^I(g)=0$ gives $\psi_0=0$. Their sum is zero and
multiplication by two on $\widehat{\mathbb Z}$ is injective;
hence all three coordinates vanish and $t=0$.
This proves \eqref{eq:cubic-normalizer-13} in
the modular quotient.

Finally both sides of \eqref{eq:cubic-normalizer-13} have exponent sum zero in the
signed braid group: $A_g\in F'$ and $f_g$ is a commutator
word. The kernel of the modular quotient is generated
topologically by $(uv)^3$, whose exponent sum is six.
Multiplication by six on $\widehat{\mathbb Z}$ is injective,
so the possible central discrepancy is zero. This proves
the signed identity exactly.
\end{proof}

In the modular quotient the identity also gives the explicit
normalized cubic relation
\[
 f_g(a,b)=f_g(r,b)f_g(r,a)^{-1}.
\]
Indeed $f_g(a,b)=\operatorname{Ad}(r)(A_g)A_g^{-1}$ and
$\operatorname{Ad}(r)(a)=b$. Substituting
$A_g^{-1}=f_g(a,r)$ and using the two-cycle identity proves
the displayed formula. Its inverse is equation~(3.10) of
\cite[p.~667]{Schneps2006Special}; that source derives it on
$\widehat{\mathrm{GS}}^1_{0,0}$ by forgetting two strands.


\begin{corollary}\label{cor:cubic-defect-zero-normalizer}
On $D=\ker\lambda\cap\bigcap_q\ker\rho_q^I$ the finite
cubic comparison defect $\epsilon_c=\tau+\eta$ vanishes.
Together with the harmonic comparison, this gives
$\chi_3=0$ on full GT.
\end{corollary}
\begin{proof}
Reduce \eqref{eq:cubic-normalizer-13} modulo eight under the signed representation.
Writing $X=U^2$ and $R=-UL$, its right side is $f_g(X,R)$:
the central sign in $q(uv)=-R$ has no effect on a commutator
word. The previously proved finite interchange map sends
$q_8(A_g)^{-1}$ to $(-I)^\tau$ and $f_g(X,R)$ to
$f_g(U,S)^{-1}=(-I)^\eta$. Hence $\tau=\eta$.
The relative extension of $\epsilon_c$ is zero on the Galois
subgroup, and $\widehat{\mathrm{GT}}=D G_{\mathbb Q}$,
so it is zero everywhere. The already proved harmonic
equality $w=\eta$ now gives $\chi_3=w+\tau=0$.
\end{proof}
